%!TeX program = xelatex
\documentclass{SYSUReport}
\usepackage{tabularx} % 在导言区添加此行
\usepackage{array}
\usepackage{booktabs}
\usepackage{caption}
\usepackage{multirow}
\usepackage{graphicx}
% 用于从 CSV 读取表格与简单绘图
\usepackage{pgfplotstable}
\usepackage{pgfplots}
\pgfplotsset{compat=1.17}
\usepackage{csvsimple}
\usepackage{siunitx}
\sisetup{detect-all=true,output-decimal-marker=.}
% 根据个人情况修改
\headl{}
\headc{}
\headr{并行程序设计与算法实验}
\lessonTitle{并行程序设计与算法实验}
\reportTitle{Lab1-基于MPI的并行矩阵乘法}
\stuname{林宏宇}
\stuid{23320093}
\inst{计算机学院}
\major{计算机科学与技术}
\date{2025年3月26日}

\begin{document}

% =============================================
% Part 1: 封面
% =============================================
\cover
\thispagestyle{empty} % 首页不显示页码
\clearpage

% % =============================================
% % Part 4： 正文内容
% % =============================================
% % 重置页码，并使用阿拉伯数字
% \pagenumbering{arabic}
% \setcounter{page}{1}

%%可选择这里也放一个标题
%\begin{center}
%    \title{ \Huge \textbf{{标题}}}
%\end{center}

\section{实验目的}
\begin{itemize}
   \item 掌握 MPI 程序的编译和运行方法。
    \item 理解 MPI 点对点通信的基本原理。
    \item 了解 MPI 程序的 GDB 调试流程。
\end{itemize}

\section{实验内容}
\begin{itemize}
   \item 使用 MPI 点对点通信实现并行矩阵乘法。
    \item 设置进程数量（1$\sim$16）及矩阵规模（128$\sim$2048）。
    \item 根据运行时间，分析程序的并行性能。
\end{itemize}


\section{实验结果}
\subsection{性能分析}
根据运行结果，填入下表以记录不同进程数和矩阵规模下的运行时间：
% Please add the following required packages to your document preamble:
% \usepackage{booktabs}
% \usepackage{multirow}
% \usepackage{graphicx}
\begin{table}[htpb]
\centering
\begin{tabular}{|c|c|c|c|c|c|}
\hline
进程数 & 128 & 256 & 512 & 1024 & 2048 \\ \hline
1      & 0.005997 s & 0.024128 s & 0.444456 s & 3.531901 s & 48.244504 s \\\hline
2      & 0.003672 s & 0.013877 s & 0.211537 s & 1.776948 s & 27.174872 s \\\hline
4      & 0.001991 s & 0.009174 s & 0.133641 s & 0.912547 s & 18.904544 s \\\hline
9      & 0.000955 s & 0.004348 s & 0.081890 s & 0.489648 s & 8.344542 s \\\hline
16     & 0.000677 s & 0.002918 s & 0.053049 s & 0.319259 s & 6.033024 s \\\hline
\end{tabular}
\caption{不同进程数和矩阵规模下的运行时间}
\label{tab:performance}
\end{table}

\subsection{结果分析}
本节使用常用并行性能指标：加速比 $S_p=T_1/T_p$ 和并行效率 $E_p=S_p/p$，其中 $T_p$ 表示 $p$ 个进程下的运行时间。基于 C 实现的完整基准结果，得到以下定量观察：

\begin{itemize}
    \item 在不同矩阵规模下的结果：
    \begin{itemize}
        \item $N=128$：$T_1=0.005997\,$s，$T_{16}=0.000677\,$s，$S_{16}\approx8.86$，$E_{16}\approx0.55$. 
        \item $N=512$：$T_1=0.444456\,$s，$T_{16}=0.053049\,$s，$S_{16}\approx8.38$，$E_{16}\approx0.52$. 
        \item $N=2048$：$T_1=48.244504\,$s，$T_{16}=6.033024\,$s，$S_{16}\approx8.00$，$E_{16}\approx0.50$. 
    \end{itemize}
    \item 结论要点：
    \begin{itemize}
        \item 随着矩阵规模增大，计算/通信比增大，整体并行效率有明显提升，说明大规模问题更能掩盖通信开销。
        \item 对于中等规模（如 $N=512,1024$），仍能观察到较好的扩展性；而对于非常小的问题（如 $N=128$），通信与调度开销占比高，效率波动较大。
        \item 未能得到理想的线性加速（$S_p=p$），主要受通信、内存带宽和负载划分影响。
    \end{itemize}
\end{itemize}


\section{讨论题}

\subsection{讨论题答案}

\begin{itemize}
\item 在内存受限情况下，如何进行大规模矩阵乘法计算？

回答：采用分块（block）方法与外存/流式计算相结合。将矩阵按行或列分块，逐块读取并进行局部乘加，避免一次性载入整个矩阵；使用内存映射（mmap）或将数据分片存储在并行文件系统上；在分布式环境下可采用行/列分片并利用 MPI 分发与汇总，必要时使用压缩表示或低秩近似减少内存占用。

\item 如何提高大规模稀疏矩阵乘法性能？

回答：使用稀疏矩阵专用存储和算法（CSR/CSC/COO）；按非零模式划分任务以减少通信与负载不均；利用缓存友好的分块与集合并（blocking）、并行化（多线程 + 分布式）；在乘法中优先跳过零值，使用稀疏 BLAS 或图计算库（例如 SuiteSparse、Intel MKL Sparse、PETSc）；对于非常大的稀疏问题，可尝试近似算法或图划分来减少跨进程通信。

\end{itemize}

\subsection{讨论题实验}

\subsubsection {分块（外存）实验}
\begin{table}[htbp]
\centering
\label{tab:mem-full}
\begin{tabular}{llrrr}
\toprule
模式 & $N$ & block & 平均时间(s) & 峰值内存(MB) \\
\midrule
naive & 512 & 64 & 0.471577 & 22.95 \\
blocked & 512 & 64 & 0.076163 & 24.44 \\
naive & 512 & 128 & 0.462930 & 23.64 \\
blocked & 512 & 128 & 0.078401 & 22.81 \\
naive & 1024 & 64 & 3.400359 & 46.24 \\
blocked & 1024 & 64 & 0.621179 & 48.40 \\
naive & 1024 & 128 & 3.434431 & 43.89 \\
blocked & 1024 & 128 & 0.618176 & 46.93 \\
\bottomrule
\end{tabular}
\caption{分块（外存）实验结果}
\end{table}

\newpage

\paragraph{分块方法对比}

下表为对比表，计算了每组配置下的加速比（naive/blocked）以及峰值内存差值（blocked - naive）：
\begin{table}[htbp]
\centering
\label{tab:mem-compare}
\begin{tabular}{rrr r r r}
\toprule
$N$ & block & naive(s) & blocked(s) & speedup & mem diff (MB) \\
\midrule
512 & 64  & 0.471577 & 0.076163 & 6.19 & +1.49 \\
512 & 128 & 0.462930 & 0.078401 & 5.91 & -0.83 \\
1024 & 64  & 3.400359 & 0.621179 & 5.48 & +2.16 \\
1024 & 128 & 3.434431 & 0.618176 & 5.56 & +3.04 \\
\bottomrule
\end{tabular}
\caption{block方法与naive方法对比（速度比 = naive / blocked）}
\end{table}

\paragraph{分块实验分析}

\begin{itemize}
    \item 在给定的四组配置中，分块实现对比天真实现的速度提升在约 $5.48\sim6.19\times$，平均约 $5.8\times$，说明分块策略显著提高了缓存局部性并降低了内存带宽压力。
    \item 峰值内存差异在数 MB 范围（从约 -0.83 MB 到 +3.04 MB 不等），表明分块实现有时会增加临时缓冲开销但总体内存变化有限。
    \item 对比不同块大小（64 vs 128），对性能影响有限：在部分配置中 128 略优，实际最优块大小仍需根据硬件缓存层次调优。
\end{itemize}

\subsubsection{稀疏矩阵实验}
\begin{table}[htbp]
\centering
\begin{tabular}{rrrl}
\toprule
$N$ & density & 平均时间(s) & 峰值内存(MB) \\
\midrule
1000 & 0.01 & 0.006954 & 21.45 \\
1000 & 0.001 & 0.001011 & 22.96 \\
2000 & 0.01 & 0.064032 & 46.88 \\
2000 & 0.001 & 0.008670 & 45.41 \\
\bottomrule
\end{tabular}
\caption{稀疏矩阵实验结果}
\end{table}

\newpage
\paragraph{稀疏实验分析}
\begin{itemize}
    \item 对于 $N=1000$，密度从 $0.001$ 增至 $0.01$，时间从 $0.001011\,$s 提升到 $0.006954\,$s，约 $6.9\times$ 更慢；对于 $N=2000$，对应比值约 $7.4\times$。这说明计算量近似随非零元素数量线性增长。
    \item 峰值内存随问题规模 $N$ 明显增加（$N=1000$ 时约 $21\sim23\,$MB，$N=2000$ 时约 $45\sim47\,$MB），对密度影响较小，表明实现的峰值内存主要受矩阵维度支配。
    \item 总结：当矩阵高度稀疏（低密度）时，稀疏专用存储与算法能显著降低计算时间；当密度增大，成本近似按非零元素数量线性增长。
\end{itemize}

\end{document}